\documentclass[12pt]{article}
\usepackage[utf8]{inputenc}
\usepackage[russian]{babel}
\usepackage{amsmath}
\usepackage{amssymb}
\usepackage{geometry}
\geometry{a4paper, margin=2cm}

\begin{document}

\section*{Задача 3.1.4}

\textbf{Условие:} Вывести двоичное представление десятичного числа. Least significant bit слева (0.5 балла).\\
in: x (1 регистр) \\
out: y (log2(x) регистров)

\section*{Решение}

\begin{verbatim}
Read 1  ; x

Load 1
JZero END_0     ; x == 0 => Out: 0

Load =3
Store 2         ; r2 := 3 - reference to the [last + 1] cell of buffer

WHILE: Load =0
    Sub 1
    JGTZ WRITE_RESULT
    JZero WRITE_RESULT  ; exit if x <= 0
    Load 1
    Div =2
    Mult =2
    Sub 1
    Mult =-1    ; r0 := x mod 2 = ((x // 2) * 2 - x) * -1
    Store *2    ; buffer[r2] := r0
    Load =1
    Add 2
    Store 2     ; r2 += 1
    Load 1
    Div =2
    Store 1     ; x := x // 2
    Jump WHILE
WRITE_RESULT: Load =3
    Sub 2
    JZero END   ; r2 == 3, all the buffer was read
    Load 2
    Sub =1
    Store 2     ; r2 -= 1
    Write *2    ; OUT <- buffer[last]
    Jump WRITE_RESULT
END_0: Write =0
END: Halt 0
\end{verbatim}

\section*{Сложность}

Худшим случаем будет x > 0, тогда цикл будет выполняться столько раз, сколько бит в двоичном представлении x. Пусть $n := \lfloor \log_2 x \rfloor + 1$.

\begin{enumerate}
    \item Инициализация выполняется за 5 шагов
    \item Первый цикл WHILE выполняется n раз
    \item Каждая успешная итерация выполняется за 17 шагов
    \item После n-й итерации выполняется ещё одна проверка условия - 4 шага
    Итого на цикл \(17n + 4\) шагов
    \item Второй цикл WRITE\_RESULT также выполняется n раз
    \item Каждая успешная итерация выполняется за 8 шагов
    \item После n-й итерации выполняется еще одна проверка условия - 3 шага
    Итого на этот цикл \(8n + 3\) шагов
    \item Halt 0 - 1 шаг
\end{enumerate}

Значит сложность в худшем случае:
\[
    T_{\max} = 5 + (17n + 4) + (8n + 3) + 1 = 25n + 13
\]
где $n = \lfloor \log_2 x \rfloor + 1$

\end{document}